\section{Conclusion}

Vibe Shuffle studies baseline-relative music selection from camera and cardiac
observations. It tracks recent movement in Valence--Arousal space, averages the
trajectory, maps that point to one of four regions, and chooses the nearest
eligible song. The trial record identifies the target, region, and distance
used for each recommendation.

In the 15-participant crossover, average mood fit was $0.56$ points higher under
Vibe (95\% CI $[-0.21,1.33]$), and liking was $0.09$ points higher (95\% CI
$[-0.52,0.71]$). The current-mood reports were subjective ratings, and the
research system produced no objective emotion label. The stored target,
signal-source summary, region, and song distance document the final decision;
the trial record omits point-level trajectory samples.

A larger preregistered study should validate the sensor mappings, use
time-weighted fusion, retain complete trial records, and include signal
ablations. If these studies establish a reproducible benefit, a later platform
study could evaluate the policy as a recommendation layer within a
music-streaming service. That evaluation must determine whether
baseline-relative selection improves the perceived fit of what plays next.
